\documentclass[10pt,a4paper]{amsart}
\usepackage[margin=19mm]{geometry}
\usepackage{amsmath,amssymb,mathtools}
\usepackage[scr=rsfs,cal=euler]{mathalfa}
\usepackage{xfrac}
\usepackage[unicode,hidelinks,bookmarks=false]{hyperref}
\newcommand{\ie}{i.e.\ }
\newcommand{\cf}{cf.\ }
\newcommand{\resp}{respectively\ }
\newcommand{\Subsection}{\subsection}
\providecommand{\nobreakdash}{\nobreak\hbox{-}}
\setlength{\emergencystretch}{2em}
\multlinegap0pt
\usepackage{amssymb}
\usepackage{mathtools}
\newtheorem{lemma}{Lemma}
\newtheorem{proposition}{Proposition}
\newcommand{\R}{\mathbb{R}}
\newcommand{\Z}{\mathbb{Z}}
\title{On the paucity of lattice triangles}
\author{David Kurniadi Angdinata, Evan Chen, Ken Ono, Jesse Thorner, Jiaxin Zhang, Jujian Zhang}
\date{Source: arXiv:2603.23928v2 (CC BY 4.0)}
\begin{document}
\maketitle

\noindent\emph{Source and licence.} On the paucity of lattice triangles. Version arXiv:2603.23928v2 (CC BY 4.0), distributed by arXiv under the Creative Commons Attribution 4.0 International licence, http://creativecommons.org/licenses/by/4.0/, as recorded in the arXiv licence metadata for this exact version and shown on the abstract page https://arxiv.org/abs/2603.23928v2. Full licence notice: \texttt{licenses/arxiv-2603.23928-CC-BY-4.0.txt}.\par\smallskip

\noindent\textit{Editorial scope.} Counted unit: the article's proposition prop:reduction (source0.tex lines 465--468). The article's complete original proof of it is delivered with it (source0.tex lines 470--513, one \texttt{proof} environment of 44 lines). No mathematics is written, repaired or re-derived: the statement and the proof are the article's own lines, in the article's own notation, and no retained line is substituted or re-flowed.  Retained uncounted as the source-local context the proof needs: lines 267--270; lines 331--334; lines 371--374; lines 380--386; lines 429--432.  Nothing was omitted from any retained span, so the package carries no omission record.  No retained line contains a citation, so the wrapper carries no bibliography and no bibliography entry.  No retained line contains a figure environment, an image-inclusion command or drawing code, so no image is delivered and the package declares no asset dependency.  Editorial formatting only, in addition: an AMS \texttt{amsart} preamble that restates the article's own declarations and macros used by the retained lines, namely the environments \texttt{lemma}, \texttt{proposition} on separate counters, together with the standard packages those lines need.  The retained environments print in this document as Lemma 1, Proposition 1 in order of appearance, whereas the article prints the same items with the numbering of its own full document.  The complete native source of the article as distributed in the arXiv e-print for this exact version is bundled unchanged as \texttt{sources/197120/source0.tex}.
\begin{lemma}
\label{lem:harmonic}
If $n\ge 2$, then $\sum_{k=1}^{n-1}w_n(k)\le \log n$.
\end{lemma}

\begin{equation}
\label{eqn:hsum}
1\le p\le h_p,\qquad 1\le q\le h_q,\qquad h_p+h_q<n.
\end{equation}

\begin{equation}
\label{eqn:detLambda}
\det\Lambda_d=[\Z^2:\Lambda_d]=d.
\end{equation}

\begin{lemma}
\label{lem:majd}
If $n\ge5$, $(p,q)\in\mathcal{H}_n$, and $d\in\mathcal{D}_n$, then
\[
\Big|N_d-\frac{H^{(d)}_pH^{(d)}_q}{d}\Big|\le\mathcal{W}_d, \qquad\text{and}\qquad |E(p,q)|\le 2^{\omega(n)+1}+\sum_{d\in\mathcal{D}_n}\mathcal{W}_d.
\]
\end{lemma}

\begin{lemma}
\label{lem:lattice}
Let $A,B>0$ and $R\coloneq \{(x,y)\colon|x|\le A,\ |y|\le B\}$.  Let $\Lambda\subseteq\R^2$ be a lattice of determinant $D$ and  $J\coloneq \#\big((\Lambda\cap R)\smallsetminus\{0\}\big)$.  If $J\ge1$, then either $J<24AB/D$, or there exists $w=(w_1,w_2)\in\Lambda\smallsetminus\{0\}$ such that $|w_1|\le4A/J$ and $|w_2|\le4B/J$.
\end{lemma}

\begin{proposition}
\label{prop:reduction}
If $(p,q)\in\mathcal{H}_n$, $M(p,q)\ge1000 \mathcal{Q}$, and $|E(p,q)|\ge M(p,q)/2$, then $(p,q)$ is resonant or degenerate.
\end{proposition}

\begin{proof}
Write $M\coloneq M(p,q)$, and suppose that $(p,q)$ is neither resonant nor degenerate.  Since $M\ge1000\cdot2^{\omega(n)}$, we have $2^{\omega(n)+1}\le M/4$.  Lemma~\ref{lem:majd} therefore gives
\begin{equation}
\label{eqn:needbig}
\sum_{d\in\mathcal{D}_n}\mathcal{W}_d \ge \frac M2-2^{\omega(n)+1} \ge \frac M4 .
\end{equation}

For each $d\in\mathcal{D}_n$, we partition the index set of $\mathcal{W}_d$ into the families
\[
\mathcal{A}_d\coloneq \{(0,\ell)\in\Lambda_d\cap(\mathcal{F}_d\times\mathcal{F}_d)\colon \ell\ne0\},\qquad \mathcal{B}_d\coloneq \{(k,0)\in\Lambda_d\cap(\mathcal{F}_d\times\mathcal{F}_d)\colon k\ne0\},
\]
and
\[
\Lambda_d(K,L)\coloneq \{(k,\ell)\in\Lambda_d\cap(\mathcal{F}_d\times \mathcal{F}_d)\colon K\le|k|<2K,\ L\le|\ell|<2L\},
\]
where $K$ and $L$ run over the powers of $2$ that are at most $d/2$.  This is indeed a partition: every $(k,\ell)$ in the index set with $k\ell\ne0$ satisfies $1\le|k|,|\ell|\le d/2$, and hence lies in exactly one of the sets $\Lambda_d(K,L)$.  The \emph{contribution} of a family $\mathcal{G}$ is $\frac1d\sum_{(k,\ell)\in\mathcal{G}}u_d(k)v_d(\ell)$, so that $\mathcal{W}_d$ equals the sum of the contributions of its families.  The total number of families is at most $2^{\omega(n)}\big(2+(\log_2n)^2\big)$, which is at most $5 \mathcal{Q}/(1+\log n)^2$.  By \eqref{eqn:needbig} and the pigeonhole principle, there exists a family with contribution at least $(1+\log n)^2M/(20 \mathcal{Q})\geq M/(20 \mathcal{Q})$.

Suppose first that this family is $\mathcal{A}_d$ for some $d$, and let $g_q\coloneq \gcd(q,d)$.  If $(0,\ell)\in\mathcal{A}_d$, then $\ell q\equiv0\pmod d$, hence $(d/g_q)\mid\ell$.  Writing $\ell=jd/g_q$ with $j\in\mathcal{F}_{g_q}\smallsetminus\{0\}$, we have $v_d(\ell)=g_q/(2|j|)$.  In particular, $\mathcal{A}_d$ is nonempty, so $g_q\ge2$.  Since $u_d(0)=H^{(d)}_p\le h_pd/n$, Lemma~\ref{lem:harmonic} gives
\[
\frac{1}{d}\sum_{(0,\ell)\in\mathcal{A}_d}u_d(0)v_d(\ell) \le \frac{h_p}{n}\sum_{j\in\mathcal{F}_{g_q}\smallsetminus\{0\}}\frac{g_q}{2|j|} \le \frac{h_p}{n} g_q\log g_q .
\]
The left-hand side is at least $(1+\log n)^2M/(20 \mathcal{Q})$.  Since $h_p<n$ by \eqref{eqn:hsum} and $\log g_q\le1+\log n$, it follows that $g_q\ge(1+\log n)M/(20 \mathcal{Q})\ge M/(1000 \mathcal{Q})$.  As $g_q$ divides $\gcd(q,n)$, the pair $(p,q)$ is degenerate, which is a contradiction.  The family $\mathcal{B}_d$ is treated symmetrically, with $\gcd(p,d)$ in place of $g_q$.

Suppose instead that this family is $\Lambda_d(K,L)$ for some $d$, $K$, and $L$.  Let $J\coloneq \#\Lambda_d(K,L)$, so that $J\ge1$.  On this family, $u_d(k)\le d/(2K)$ and $v_d(\ell)\le d/(2L)$, so the contribution is at most $dJ/(4KL)$.  It follows that $dJ/(4KL)\ge M/(20 \mathcal{Q})$; that is,
\begin{equation}
\label{eqn:KLbound}
KL\ \le\ \frac{5 d J \mathcal{Q}}{M}.
\end{equation}
We now apply Lemma~\ref{lem:lattice} with $\Lambda=\Lambda_d$, $D=d$ (see \eqref{eqn:detLambda}), $A=2K$, and $B=2L$.  Let $J'$ denote the quantity denoted $J$ in that lemma.  The box $R$ contains $\Lambda_d(K,L)$, so $J'\ge J\ge1$.

If the first alternative of Lemma~\ref{lem:lattice} holds, then $J\le J'<24AB/D=96KL/d$.  Combined with \eqref{eqn:KLbound}, this yields $M<480 \mathcal{Q}$, which contradicts $M\ge1000 \mathcal{Q}$.  Hence the second alternative holds: since $4A/J'\le8K/J$ and $4B/J'\le8L/J$, there is $(k_0,\ell_0)\in\Lambda_d\smallsetminus\{0\}$ with $|k_0|\le8K/J$ and $|\ell_0|\le8L/J$.  Let $(k_0',\ell_0')\in\mathcal{F}_d\times\mathcal{F}_d$ be the representative of $(k_0,\ell_0)$ modulo $d$, so that $(k_0',\ell_0')\in\Lambda_d$, $|k_0'|\le|k_0|$, and $|\ell_0'|\le|\ell_0|$.  Three cases arise.

First, suppose that $(k_0',\ell_0')=(0,0)$.  Then $d$ divides both $k_0$ and $\ell_0$, so at least one of $k_0$ and $\ell_0$ is a nonzero multiple of $d$; say $|\ell_0|\ge d$, the case $|k_0|\ge d$ being symmetric.  Then $8L/J\ge|\ell_0|\ge d$, so $KL\ge L\ge Jd/8$.  Combined with \eqref{eqn:KLbound}, this yields $M\le40 \mathcal{Q}$, which is a contradiction.

Second, suppose that $k_0'\ell_0'\ne0$.  By \eqref{eqn:KLbound},
\[
|k_0'\ell_0'|\le\frac{64 KL}{J^2}\le\frac{320 d \mathcal{Q}}{J M}\le\frac{1000 d \mathcal{Q}}{M},
\]
so $(p,q)$ is resonant, which is a contradiction.

Third, suppose that exactly one of $k_0'$ and $\ell_0'$ vanishes; say $\ell_0'=0$, the case $k_0'=0$ being symmetric.  Let $g_p\coloneq \gcd(p,d)$.  Then $k_0'p\equiv0\pmod d$, so $(d/g_p)\mid k_0'$, and therefore $d/g_p\le|k_0'|\le8K$.  Now fix $\ell$ with $L\le|\ell|<2L$; there are at most $2L$ such $\ell$.  The integers $k$ with $(k,\ell)\in\Lambda_d$, if any, lie in a single residue class modulo $d/g_p$, so at most $4Kg_p/d+1$ of them satisfy $|k|<2K$; since $d/g_p\le8K$, this count is at most $12Kg_p/d$.  Hence $J\le24KLg_p/d$.  The contribution of $\Lambda_d(K,L)$ is therefore at most $dJ/(4KL)\le6g_p$, while it is at least $M/(20 \mathcal{Q})$, so $g_p\ge M/(120 \mathcal{Q})\ge M/(1000 \mathcal{Q})$.  As $g_p$ divides $\gcd(p,n)$, the pair $(p,q)$ is degenerate, which is a contradiction.

These cases are exhaustive, so the proof is complete.
\end{proof}

\end{document}
